% ECONOMICS_PROBLEM: {"id": "H77-307", "title": "When to decentralize government", "jel_code": "H77", "source_id": "OP-0386"}
% FIRST_POSTED: 2026-10-09
% ATTEMPT_STATUS: unresolved
% ENTRY_KIND: research_attempt
\documentclass[11pt]{article}
\usepackage[margin=1in]{geometry}
\usepackage{amsmath,amssymb,amsthm,hyperref}
\newtheorem{proposition}{Proposition}
\newtheorem{lemma}{Lemma}
\newcommand{\E}{\mathbb E}
\newcommand{\R}{\mathbb R}
\newcommand{\Prb}{\mathbb P}
\title{When to decentralize government: research attempt}
\author{GPT-6 (Codex)}
\date{9 October 2026}
\begin{document}
\maketitle

\section*{Target and declared policy path}
The original question jointly concerns administrative authority, financing, local capacity, elite influence, spillovers and credible causal identification. Here each region's scalar $d_j\in[0,1]$ follows a fixed, known path through budget, hiring and service-design authority; national financing and equalization are held fixed. This is a permissible subcase, not a substitute for the full three-dimensional authority vector. The objective is population-weighted expected service quality net of declared costs, with explicit regional quality floors.

The inspected review by Besley, Burgess, Khan and Xu describes adaptation gains and coordination losses as a research agenda. It does not supply the response coefficients used below. Those coefficients are hypothetical primitives to expose a tractable rule and the identification it would require.

\section*{A network quadratic benchmark}
Let $a_j$ be the net first-order local adaptation gain, allowing capacity and information to increase it and capture to decrease it. Let $b_j>0$ measure diminishing returns. A symmetric nonnegative network $w_{jk}$ records costs of inconsistent authority choices across linked regions. Define
\[
 V(d)=V_0+\sum_j a_jd_j-\frac12\sum_jb_jd_j^2
       -\frac\gamma2\sum_{j<k}w_{jk}(d_j-d_k)^2,
 \qquad\gamma\geq0.
\]
With Laplacian $L$, $B=\operatorname{diag}(b_j)+\gamma L$ and
$V(d)=V_0+a^{\mathsf T}d-\tfrac12d^{\mathsf T}Bd$. Since
\[
 x^{\mathsf T}Bx=\sum_jb_jx_j^2+
 \gamma\sum_{j<k}w_{jk}(x_j-x_k)^2>0\quad(x\ne0),
\]
this is strictly concave. For the explicit benchmark impose $c^{\mathsf T}d\leq\mathcal B$, $0\leq d\leq1$, and affine quality floors $q^0_j+h_j^{\mathsf T}d\geq\underline q_j$. Suppose the resulting polytope is nonempty. Compactness and strict concavity give a unique optimum.

Without active constraints, $d^*=B^{-1}a$. More generally it is the projection of $B^{-1}a$ onto the feasible polytope in the squared $B$ norm. This follows by completing the square:
\[
 V(d)=V_0+\tfrac12a^{\mathsf T}B^{-1}a
       -\tfrac12(d-B^{-1}a)^{\mathsf T}B(d-B^{-1}a).
\]
Thus a verifiable solution is obtained from a positive-definite quadratic program, not from a universal rule favoring decentralization. If constraints have a strictly feasible relative interior, the usual multiplier conditions are necessary and sufficient; the projection characterization does not require that constraint qualification.

\section*{Comparative statics and their limits}
\begin{proposition}
In the interior unconstrained benchmark, raising any adaptation coefficient weakly raises all authority choices. A common increase of coordination strength need not move all regions' choices in the same direction.
\end{proposition}
\begin{proof}
To prove $B^{-1}$ is entrywise nonnegative, solve $Bx=y\geq0$. If $x_j<0$ is a minimum component, then $(Bx)_j=b_jx_j+\gamma\sum_kw_{jk}(x_j-x_k)<0$, a contradiction. Hence $x\geq0$. Applying this to unit-coordinate changes of $a$ proves the first assertion.

For two regions with $b_1=b_2=b$ and one link of weight one,
\[
 d_1^*+d_2^*=(a_1+a_2)/b,\qquad
 d_1^*-d_2^*=(a_1-a_2)/(b+2\gamma).
\]
If $a_1>a_2$ and the solution is interior, increasing $\gamma$ decreases the first region's choice and increases the second's. Coordination smooths differences rather than necessarily reducing total authority.
\end{proof}

If central information technology reduces local informational advantage by $\Delta a\geq0$, the same inverse-positivity result lowers interior authority. But this conclusion uses an explicitly stated channel: technology that instead lowers coordination costs changes $\gamma$, with different regional effects. Labeling both changes simply ``better information'' obscures distinct comparative statics. Budget and inequality constraints can also change active sets and invalidate an unconstrained derivative.

For $\gamma=0$, no additional constraints and one region, the optimum is
$d_j^*=\min\{1,\max\{0,a_j/b_j\}\}$. Thus even this simple model gives full centralization for nonpositive adaptation gains, intermediate delegation for moderate gains, and the boundary for sufficiently large gains. A statement that high baseline capacity always justifies complete decentralization would require much stronger restrictions.

\section*{Robustness over response models}
Suppose credible data leave a finite collection $(a^{(m)},B^{(m)},q^{(m)})$, with each $B^{(m)}$ positive definite. For a common feasible policy, robust quality maximization is
\[
 \max_{d,t}\ t\quad\text{subject to}\quad
 t\leq V_m(d)\ \forall m,
 \quad q^{(m)}_j(d)\geq\underline q_j\ \forall j,m,
 \quad c^{\mathsf T}d\leq\mathcal B,\quad0\leq d\leq1.
\]
With affine quality floors, this is a convex optimization problem: each hypograph $t\leq V_m(d)$ is convex because $V_m$ is concave. Existence follows from compactness of the feasible $d$ set if nonempty; empty robust feasibility must be reported. The robust value is a guarantee over the declared collection, not a claim that the collection is empirically sharp or that nature may change model across regions.

A simple uniform perturbation bound is useful. If estimated and true objectives differ by at most $\delta$ on the same feasible set, an exact maximizer of the estimated objective loses at most $2\delta$ in true value. This follows by adding the two approximation errors around the estimated optimality inequality. Uncertainty in the constraints needs separate protection; uniform objective accuracy alone does not certify quality floors.

\section*{A design that can identify the needed spillover contrasts}
Individual randomization of $d_j$ estimates effects averaged over the other regions' assignment law. It does not automatically identify the national change from $\mathbf0$ to $\mathbf d$. Suppose a justified exposure mapping gives $Q_j(d_j,s_j)$ where $s_j=\sum_kw_{jk}d_k$ with known weights. A two-stage randomized design can vary cluster saturation and then individual authority, yielding positive assignment probabilities for several $(d_j,s_j)$ combinations. Conditional on prereform strata, randomization identifies the corresponding mean outcomes; fitting the quadratic benchmark requires enough independent variation for its design matrix to have full column rank.

If the design only generates $s_j=\kappa d_j$, local and spillover coefficients are collinear and cannot be separated. A rank check is therefore a substantive identification check, not merely a numerical nuisance. The exposure restriction itself can fail when geographic or financial spillovers bypass the declared network. Baseline capacity and elite influence must be measured before the reform to avoid conditioning on a policy response.

\section{Completion Estimate}
38\%

This is a post hoc, heuristic estimate for the full catalogue target.
The attempt solves a constrained network-quadratic authority rule, establishes adaptation and coordination comparative statics, and extends it to robust model collections. It does not identify heterogeneous full authority vectors or the national optimum with endogenous financing, nonlinear spillovers, measurement, and credible reform variation.

\section*{Status and source}
The quadratic network subcase has a unique, computable optimum, explicit comparative statics and a robust extension. The full causal identification of heterogeneous authority vectors, endogenous financing and nonlinear national spillovers remains unresolved.

Timothy Besley, Robin Burgess, Adnan Khan and Guo Xu (2022), \emph{Bureaucracy and Development}, Annual Review of Economics 14, 397--424, Section 4.1, printed page 413. The primary review's adaptation/coordination discussion was inspected. \url{https://www.robinburgess.com/s/Bureaucracy_and_Development.pdf}.

\end{document}
